\section{Results}\label{sec:results}

\subsection{Mechanistic diagnosis}\label{sec:mechanism}

The deployed inverse document frequency uses the vocabulary size,
$\mathrm{idf}_{\mathrm{dep}}(v)=\log_{10}(K/\mathrm{df}(v))$, whereas canonical
TF-IDF uses the number of documents,
$\mathrm{idf}_{\mathrm{can}}(v)=\log_{10}(N/\mathrm{df}(v))$. Their relation is
exact:
\begin{equation}
S_{\mathrm{dep}}(s)=S_{\mathrm{can}}(s)+D_s\log_{10}(K/N),
\label{eq:idfrelation}
\end{equation}
where $D_s$ is the number of descriptors. The formulations therefore differ by
an additive count term, not by rescaling. The identity holds within
$3\times10^{-11}$ on all 93 segment slots. At cluster scale
$N\in\{2,3,4\}$, so the additional term rewards count and contributes a median
98.3\% of the deployed score. On a database of approximately $10^4$ documents,
the term is negative and reverses the count ordering.

The more general factorization separates the two available channels:
\begin{equation}
\mathrm{score}_i=\sum_v\mathrm{tf}_{iv}\mathrm{idf}_v
=D_i\underbrace{\sum_v p_i(v)\mathrm{idf}_v}_{\overline{\mathrm{idf}}_i}.
\label{eq:factor}
\end{equation}
Thus a score can vary through how many descriptors the extractor finds and
through their mean IDF weight. The identity alone does not prove that either channel
is uninformative; the relative contribution is measured empirically.

\begin{figure}[!ht]
\centering
\begin{tikzpicture}[
  box/.style={rectangle,rounded corners=3pt,minimum height=1.15cm,
    text width=3.45cm,align=center,font=\sffamily\scriptsize,
    draw=#1,fill=#1!8,line width=.9pt},
  arr/.style={-{Stealth[length=5pt]},semithick,rulegray}]
\node[box=catbovw] (score) {\textbf{BoVW+TF-IDF score}\\$D_i\,\overline{\mathrm{idf}}_i$};
\node[box=accentorange,below left=.65cm and .55cm of score] (count)
 {\textbf{Descriptor-count channel}\\KonIQ CV .938\\$\rho(S,D)=.998$};
\node[box=accentpurple,below right=.65cm and .55cm of score] (rare)
 {\textbf{Mean-IDF-weight channel}\\KonIQ CV .098\\19-extractor mean $|\rho|=.090$};
\node[box=accentred,below=.7cm of count] (dep)
 {\textbf{Deployed IDF shift}\\$D_i\log_{10}(K/N)$\\sign changes with document count};
\node[box=accentgreen,below=.7cm of rare] (test)
 {\textbf{Controlled consequence}\\65.1\% over 800 pairs\\noise 8\%; JPEG 6\%};
\draw[arr] (score)--(count); \draw[arr] (score)--(rare);
\draw[arr] (count)--(dep); \draw[arr] (rare)--(test);
\draw[arr] (dep)--(test);
\end{tikzpicture}
\caption{Measured mechanism. The factorization identifies two channels; the
data determine their relative variation and connect count dominance to the
controlled-degradation failures.}
\label{fig:mechanism}
\end{figure}

On KonIQ-10k, the coefficient of variation is .938 for $D_i$ and .098 for mean
IDF weight. The canonical score consequently correlates .998 with descriptor
count. Its association with human scores is statistically equivalent, within
the operational $\pm.02$ margin, to that of keypoint density:
$\rho^{*}=.447$ for both, with score-to-keypoint correlation .998. Steiger's
$Z=.199$ is below the Bonferroni-corrected critical value of 3.95, and both
one-sided equivalence tests reject. KADID gives the same association-level
result under reference-block resampling: .125 for canonical BoVW and .123 for
keypoint density, paired difference $+.0024$, 95\% CI
$[-.0019,+.0084]$. This is equivalence of measured MOS/DMOS association, not
identity of the scores or constructs. The deployed formulation is almost exactly the sign-reversed
count on the large databases ($-.9999$ on KonIQ, $-.9998$ on KADID and
$-.9947$ on CLIVE). In the deployment regime its own descriptor count remains
dominant, $\rho(S_{\mathrm{dep}},D)=+.994$.

Replacing SIFT with 19 convolutional extractors does not raise the external
ceiling. Their mean $|\rho|$ is .090 on KonIQ and .112 on CLIVE, with no value
above .238 on either database; KADID ranges from .156 to .418. Fixed spatial
feature maps remove per-image count variation on the image databases, leaving
the rarity channel measured by the ablation. On broadcast segments, however,
the upstream key-frame selector varies the number of retained frames, so it
reintroduces count variation: across the 19 variants,
$\rho(S,D)=+.981$ to $+.997$. These results explain why exchanging the
extractor changes individual choices without making the ordering a better
predictor of quality. The explanation is limited to the evaluated vocabulary,
normalization and extraction regimes.

Simple normalization does not repair the mechanism. Dividing by descriptor
count exposes the weak mean-IDF-weight channel already measured by the ablation. Dividing
by frame count is also mismatched to the deployed operation because its internal
key-frame selector makes descriptor count nearly independent of duration. On the
42 clusters, frame normalization changes the uncalibrated DOVER two-branch-sum
output only from 47.4\% to 50.6\% and induces an inverse duration association
($\rho=-.52$). Full normalization and pooling results are in Tables~S6 and S9.

\subsection{External validation and screening}\label{sec:screen_results}

Within-reference KADID sensitivity separates the deployed operation from
specialized quality methods. Its mean direction-aligned correlation is $-.077$
(95\% CI $[-.124,-.030]$), with only 5 of 25 types above .50. LIQE,
QualiCLIP and CLIP-IQA$+$ reach .807, .849 and .808, respectively, each with 22
types above .50. Most remaining NR-IQA candidates also respond broadly, with
means from .607 to .733.

Table~\ref{tab:screening} applies the same complete-case rule to every method.
Released-weight exposure is visible because an in-database correlation is not a
generalization result. The generalization column reports the lowest signed
$\rho^{*}$ over every unseen database, rather than using an absolute value that
could hide inversion.

\begin{table}[!ht]
\centering
\caption{Complete-case method screening. Sensitivity reports the within-reference
mean and number of KADID types above .50. Generalization is the lowest signed
$\rho^{*}$ over all unseen databases. Floor is normalized agreement with Q-Align
on the common 42 clusters and its topic-block lower bound (L). Sup.: K =
KonIQ-10k opinion training; L = LIVE; OU = opinion-unaware; -- = no opinion
training; M = multi-database mixture. Q-Align is a reference, not a candidate.}
\label{tab:screening}
\scriptsize
\setlength{\tabcolsep}{1.7pt}
\renewcommand{\arraystretch}{1.08}
\begin{tabular}{@{}lcccccc@{}}
\toprule
\rowcolor{headerblue}
\hd{Method} & \hd{Sup.} & \hd{Pol.} & \hd{Sensitivity} &
\hd{Lowest unseen $\rho^{*}$} & \hd{Floor; L} & \hd{Eligible} \\
\midrule
\rowcolor{accentgreen!15}\textbf{LIQE} & K & \polup & .807; 22 & .723 & .551; .284 & \textbf{yes} \\
\rowcolor{accentgreen!15}\textbf{QualiCLIP} & OU & \polup & .849; 22 & .665 & .371; .143 & \textbf{yes} \\
\rowcolor{accentgreen!15}\textbf{CLIP-IQA$+$} & K & \polup & .808; 22 & .654 & .236; .034 & \textbf{yes} \\
\midrule
TOPIQ & K & \polup & .721; 20 & .511 & .506; .289 & no \\
\rowcolor{rowlight}MUSIQ & K & \polup & .733; 20 & .556 & .551; .297 & no \\
DBCNN & K & \polup & .668; 19 & .484 & .461; .274 & no \\
\rowcolor{rowlight}HyperIQA & K & \polup & .673; 19 & .462 & .371; .122 & no \\
NIQE & OU & \poldown & .607; 19 & .368 & .281; .012 & no \\
\rowcolor{rowlight}BRISQUE & L & \poldown & .657; 19 & .236 & .371; .136 & no \\
CLIP-IQA & OU & \polup & .645; 20 & .501 & .371; .045 & no \\
\rowcolor{catbovw!12}SIFT+BoVW (deployed) & -- & \polup & $-.077$; 5 & $-.445$ & .056; $-.308$ & no \\
\midrule
\rowcolor{rowlight}Q-Align & M & \polup & .880; 23 & -- & -- & reference \\
\bottomrule
\end{tabular}
\end{table}

Three methods are eligible: LIQE, QualiCLIP and CLIP-IQA$+$. This is an
eligibility set, not a ranking. QualiCLIP has the largest sensitivity mean,
whereas LIQE has the largest generalization margin. The set is unchanged when
the generalization threshold varies from .60 to .65 or the sensitivity coverage
cut varies from 12 to 22 types (Table~S10); at .55 MUSIQ is added, and at .70
only LIQE remains. ARNIQA and TReS are reported on completed axes only, so
neither enters this complete-case claim.

The database-level values show why the universal unseen-database clause matters.
LIQE obtains $\rho^{*}=.723$ on KADID and .865 on CLIVE; QualiCLIP obtains .665
and .757; CLIP-IQA$+$ obtains .654 and .805. TOPIQ and MUSIQ are strong on the
unseen authentic database (.811 and .789) but fall to .511 and .556 on KADID.
DBCNN and HyperIQA similarly fall from .756 and .750 on CLIVE to .484 and .462
on KADID. An existential rule would admit methods whose favourable authentic
result masks failure under controlled synthetic distortion. Conversely, the
three eligible methods pass the floor with overlapping topic-block intervals,
so the small in-domain sample cannot order them.

The floor axis illustrates that limitation directly. LIQE agrees with Q-Align
on 32 of 42 clusters ($\kappa_{\mathrm{uniform}}=.551$, 95\% CI
$[.284,.786]$); QualiCLIP agrees on 28 ($.371$, $[.143,.597]$); and
CLIP-IQA$+$ on 25 ($.236$, $[.034,.435]$). The deployed operation agrees on 21,
only 1.3 decisions above the uniform expectation, and its normalized interval
crosses chance ($.056$, $[-.308,.412]$). The candidates pass the floor, but the
broadly overlapping intervals do not support pairwise ordering. Results on the
15 human-annotated redundancy clusters are retained in Tables~S7--S8 rather than
mixed with the operational 42-cluster set.

The full P.1401 triad and database-level correlations are in Tables~S4 and S15.
The latter reports training-database contrasts descriptively, because database,
content, architecture and objective are not matched. Conditioning on detectable
keypoint density strongly attenuates the BoVW score but is not interpreted as a
causal partition of ``content density'' (Table~S16). Cost is also secondary to
eligibility: in the common environment the deployed operation takes 182~s per
cluster and LIQE 1.9~s (Table~S14).

\subsection{Constructed broadcast degradation}\label{sec:constructed_results}

Table~\ref{tab:constructed} removes the reference model from the decision: each
pair contains the same content before and after a known degradation. The
deployed operation is correct on 65.1\% of 800 pairs. LIQE reaches 99.5\%; the
paired difference is $+34.4$ percentage points, segment-block 95\% CI
$[+31.1,+37.7]$. Every evaluated NR-IQA candidate exceeds the deployed score.

\begin{table}[!ht]
\centering
\caption{Strict accuracy (\%) on 800 unmodified--degraded broadcast pairs: 20
segments, eight types and five levels. The unmodified member is the
construction-defined higher-fidelity signal; values below the 50\% chance floor
are systematic inversions. Ties remain in the denominator.}
\label{tab:constructed}
\scriptsize
\setlength{\tabcolsep}{1.9pt}
\begin{tabular}{@{}lccccccccc@{}}
\toprule
\textbf{Method} & \textbf{dark} & \textbf{motion} & \textbf{blur} & \textbf{H.264}
& \textbf{rescale} & \textbf{satur.} & \textbf{noise} & \textbf{JPEG} & \textbf{mean} \\
\midrule
\rowcolor{catbovw!12}SIFT+BoVW (deployed) & 100 & 100 & 90 & 90 & 88 &
\textcolor{accentred}{\textbf{39}} & \textcolor{accentred}{\textbf{8}} &
\textcolor{accentred}{\textbf{6}} & \textbf{65.1} \\
\midrule
LIQE & 97 & 100 & 100 & 100 & 100 & 100 & 99 & 100 & \textbf{99.5} \\
TOPIQ & 93 & 100 & 100 & 100 & 100 & 93 & 100 & 100 & 98.2 \\
MUSIQ & 93 & 100 & 100 & 100 & 100 & 86 & 100 & 100 & 97.4 \\
QualiCLIP & 99 & 100 & 95 & 100 & 96 & 76 & 96 & 71 & 91.6 \\
ARNIQA & 91 & 100 & 99 & 98 & 100 & 61 & 82 & 98 & 91.1 \\
DBCNN & 84 & 97 & 92 & 97 & 95 & 92 & 85 & 81 & 90.4 \\
NIQE & 100 & 100 & 98 & 97 & 98 & \textcolor{accentred}{\textbf{0}} & 100 & 96 & 86.1 \\
CLIP-IQA$+$ & 76 & 99 & 79 & 100 & 91 & 87 & 84 & 64 & 85.0 \\
HyperIQA & 47 & 100 & 100 & 100 & 100 & \textcolor{accentred}{\textbf{30}} & 100 & 100 & 84.6 \\
CLIP-IQA & 58 & 100 & 83 & 100 & 92 & 91 & 55 & 76 & 81.9 \\
BRISQUE & 100 & 100 & 100 & 99 & 100 & \textcolor{accentred}{\textbf{0}} & 67 & 80 & 80.8 \\
\bottomrule
\end{tabular}
\end{table}

No method ties on these pairs. Excluding darkening and desaturation leaves 600
comparisons and does not repair the deployed operation: its accuracy is 63.7\%,
whereas LIQE reaches 99.8\% and Q-Align 100\% (Table~S12). The main diagnosis
therefore does not depend on treating those two transformations as universally
worse in human perception.

The average conceals the diagnostic result. SIFT+BoVW is at or above 88\% on
five types but selects the degraded copy in 92\% of white-noise pairs and 94\%
of JPEG pairs. Within segment, its score increases with severity for noise
($\rho=+.720$, CI $[+.455,+.925]$) and JPEG ($+.640$, CI
$[+.395,+.850]$). The mechanism is directly visible in the saved descriptor
counts. The degraded copy has more descriptors in 92\% of noise pairs (mean
change $+469$) and 95\% of JPEG pairs (mean change $+289$), with mean noise
count increasing at every severity level. These measurements are consistent
with noise and blocking creating detectable extrema that increase the count
channel in Eq.~\ref{eq:factor}; structure-removing distortions instead happen
to move count in the direction of fidelity. The score is therefore responsive,
but not consistently to quality (Table~S12).

Q-Align is correct on 98.2\% of these pairs (95\% CI $[97.2,99.1]$), which
supports its technical-degradation competence on this broadcast content. It
does not validate compositional preference. The exact downstream DOVER-T output
reaches 97.6\% ($[96.1,99.0]$), or 99.8\% ($[99.5,100]$) without darkening and
desaturation; neither result is a human judgment. The 19
extractor substitutions average 44.9\% (CI $[41.0,48.9]$) and remain dominated
by their descriptor count, corroborating the mechanism while refuting the
recorded prediction that fixed feature maps would force chance performance.
The full prediction and overlay audits are in Table~S12.

\subsection{Pipeline substitution}\label{sec:e2e}

LIQE is used for the direct operation substitution because it has the largest
generalization margin among the three eligible methods and near-ceiling accuracy
on the constructed pairs. With the same ten uniformly spaced frames and mean
pooling, SIFT+BoVW and LIQE still differ in 16 of 42 clusters. The LIQE--SIFT
DOVER-T contrast is $+27.1$ range points (topic CI $[5.6,47.8]$; hierarchical
$[2.4,52.5]$) and $+.0074$ raw, with both intervals positive (Table~S11).
Table~\ref{tab:e2e} reports the complete-operation comparison and all eligible
methods; screening does not rank them.

\begin{table}[!ht]
\centering
\caption{Model-reference output over the 42 operational redundancy clusters.
Changed counts representatives differing from the deployed operation. DOVER-T
is the raw DOVER technical branch and is not used in screening; Q-Align defines
the in-domain screening floor. Both outputs use the global-achievable-range
estimand in Eq.~\ref{eq:globalrange}; Max is the percentage of clusters in which
the selected representative has the highest DOVER-T score.}
\label{tab:e2e}
\small
\setlength{\tabcolsep}{5pt}
\renewcommand{\arraystretch}{1.08}
\begin{tabular}{@{}lcrrr@{}}
\toprule
\rowcolor{headerblue}
\hd{Quality operation} & \hd{Changed} & \hd{DOVER-T} & \hd{Max} & \hd{Q-Align} \\
\midrule
\rowcolor{catbovw!12}SIFT+BoVW & --- & 50.5\% & 54.8\% & 48.6\% \\
\rowcolor{accentgreen!15}\textbf{LIQE} & \textbf{22/42} & \textbf{76.0\%} & 59.5\% & 90.2\% \\
QualiCLIP & 19/42 & 71.9\% & 59.5\% & 80.6\% \\
\rowcolor{rowlight}CLIP-IQA$+$ & 22/42 & 70.7\% & 50.0\% & 74.3\% \\
\bottomrule
\end{tabular}
\end{table}

Replacing the deployed visual-quality operation with the LIQE-based operation
changes 22 of 42 representatives. On the primary DOVER technical branch, the
global-achievable-range point estimate rises from 50.5\% to 76.0\%, a paired
increase of $+25.5$ points. Its topic-block interval is $[-12.1,+54.9]$ and its
hierarchical interval is $[-17.2,+59.6]$; neither confirms a population-level
effect. The raw branch-score difference is $+.0069$ (topic-block
$[-.0028,+.0178]$; hierarchical $[-.0041,+.0197]$), while mean per-cluster
position and selected-maximum rate increase by 7.2 and 4.8 points, respectively.

No cluster has zero DOVER-T range. Excluding the smallest 10\% or 25\% of ranges
gives global-range differences of $+25.9$ and $+28.1$ points, so the positive
point estimate is not generated solely by near-zero ranges. As a construct
sensitivity analysis, the uncalibrated sum of DOVER's technical and aesthetic
branches reproduces the earlier $+30.3$-point contrast, but that sum is neither
the official standardized fused score nor a purely technical measure
(Table~S11).

The two model references have positive point estimates for the substitution but not on
the ordering within the eligible set. That disagreement, together with
overlapping intervals, is why the study reports eligibility rather than a
winner. More importantly, neither column measures human preference for the
composed summary. The supported end-to-end conclusion is that the operation
substitution materially changes the selected evidence; the model-reference
output effects are positive but imprecise, and improved viewer experience
remains untested.
